\section{Integration with the Compiler}
\label{sec:compiler}

\subsection{A pass after type checking}

The ownership pass runs where the effect inference runs, after every type is solved and before the
interface is published. Futhark moved its alias checking out of its type checker in 2026 for the same
reason: the old design ``had always been hindered by having to work with incomplete type
information'', whereas a separate pass is ``given a fully well-typed program and simply had to figure
out whether any of the in-place constraints were violated (essentially, whether in-place updates are
semantically observable)'' \citep{henriksen2026rewrite}. The pass never touches unification or the
type store, and its messages carry solved types. It needs three things the checker does not keep today:
the solved type of every list-typed instruction, to decide paths and crossing; a cell-site identifier
per instruction (the instruction index); and the source span of every holder, for messages.

\subsection{Effects and dispatch}

The two effect bits and the ownership classes ride on the same function-type sites, are solved in the
same per-module phase, and are published in the same kind of interface block. They interact in one
place: a function that \emph{suspends} parks its locals in a continuation, which is one-shot (A7), so a
parked local is the same holder it was and liveness is unchanged. An impure call is a call like any
other for holders, and the optimiser's duty to keep impure calls in place is what A5 extends. A method
call's callee is the term the dispatch table resolves, and its summary is that term's; evidence a
function receives is a class instantiated at its caller (\S\ref{sec:inf-generic}).

\subsection{The write-set analysis}

The direct platform's whole-program write-set analysis and the ownership checker meet only at
\code{update}'s entry. The checker computes the entry state itself (\S\ref{sec:hard-tea}), because the
write-set analysis does not analyse commands and records no aliasing between two model paths. The
write-set consumer needs no new input, but one word of its soundness statement changes: its notion
of ``changed'' was ``not the same value, not \code{===}''. Under in-place writes a write inside
\code{model.rows} leaves the list \code{===} to its old self. The write set still \emph{names} the
written path, because it is read off the arm's text rather than off identity, so a consumer that runs
the DOM updates the write set names stays correct; what it may no longer do is conclude ``unchanged''
from \code{===} on a list path, which is A4. Its own requirement that an in-place list write never
touch a list the program still holds becomes the theorem of \S\ref{sec:soundness}.

\subsection{The backend}
\label{sec:compiler-backend}

Every accepted demand is in place, so the backend needs no per-site decision: the core library's
writers \emph{are} the in-place operations (Table~\ref{tab:backend}).

\begin{table}[t]
\centering\small
\begin{tabularx}{\linewidth}{@{}>{\raggedright\arraybackslash}p{0.27\linewidth}>{\raggedright\arraybackslash}p{0.3\linewidth}X@{}}
\toprule
Primitive & Today & Under the rule \\
\midrule
\code{set}, \code{update}, \code{swap} & copy up to 256 elements, else trie path copy & \code{a[i] = v} on the base; $O(1)$ \\
\code{push} & copy under 32 elements, else claim the tail & \code{a.push(v)}; amortised $O(1)$ \\
\code{pop} & copy up to 256 elements, else a new header & \code{a.pop()}; $O(1)$ \\
\code{insertAt}, \code{removeAt} & $O(n)$ plain copy & \code{a.splice(i, …)}; $O(n)$ move, as hand-written code does \\
\code{cons} (\code{[ x, …xs ]}) & claim a head slot, short copy, or convert & \code{a.unshift(x)}; $O(n)$ \\
\code{append} & push onto a trie or concatenate & push each element; $O(|ys|)$ \\
\code{copy} & --- & \code{a.slice()} \\
write through a view & --- & the same on the base at offset $o + i$; a new header for \code{push}/\code{pop} \\
\bottomrule
\end{tabularx}
\caption{List primitives today and under the rule.}
\label{tab:backend}
\end{table}

\emph{The trie goes.} Its claims, its offsets, its head and tail buffers, the cache of its plain form
and its thresholds all exist to make writes to \emph{shared} lists cheap, and under the rule no shared
list is ever written. What stays is the plain form, the view form for a pattern's \code{rest} (with the
backend's loop rewrite still removing most views), the array builder used by functions that build their
result in tail position, and the three-fact protocol for readers. Our prototype study measured the list
runtime at 706 bytes compressed without the trie, against 1\,514 with it.

The backend must still keep the identity promises the library chooses to keep (\S\ref{sec:disc-options})
and the linearity of the builder. \code{++} on lists lowers to \code{append}, a demand on its first
argument; the generic \code{++} on other appendable types stays a copy.

\subsection{The release optimiser}

The verdict is a function of the source, and \code{--release} emits the same in-place calls as a
development build, so beni's promise that a release build behaves exactly as a development build holds
with no new exception---\emph{provided the optimiser treats a demand as a barrier} (A5). With that, the
rest of the optimiser's contract carries over: an unused binding whose right-hand side is a pure call
may be dropped (an unused \code{List.set} is then not performed, which no reader can tell); two surviving
calls are never reordered; inlining never duplicates an argument's evaluation; single-use inlining may
inline a $\mathit{once}$ closure's body into its one call; specialised copies inherit the general
verdict. Renaming and minification touch names, not writes, and \code{copy} ships only when reached.

\subsection{Records in place}
\label{sec:compiler-records}

Records are outside this paper's scope, but the machinery applies: a record path is a ``cell'' whose
demand is an in-place record update, and its holders are the same places. The difference is that the
renderer compares records by identity everywhere---a reference check per dynamic hole is what makes
beni's compiled templates fast---so in-place records need a stronger form of A4 than lists do.

\subsection{Incremental builds}

The summary block is in the interface and its hash, and beni's interface firewall fires when it
changes. Two consequences need measurement: churn, since a body edit that turns a borrowed parameter
consumed re-checks every importer (necessarily, since their side conditions change); and the warm
rebuild budget, since a body edit that changes a summary \emph{is} a signature edit under this design.
